\documentclass[11pt,twocolumn]{article}

\usepackage[margin=0.85in]{geometry}
\usepackage{amsmath,amssymb,amsthm}
\usepackage{graphicx}
\usepackage{hyperref}
\usepackage{listings}
\usepackage{xcolor}
\usepackage{booktabs}
\usepackage{longtable}
\usepackage{enumitem}
\usepackage{algorithm}
\usepackage{algpseudocode}
\usepackage{multicol}
\usepackage{float}
\usepackage{caption}

\hypersetup{
  colorlinks=true,
  linkcolor=blue!60!black,
  citecolor=green!40!black,
  urlcolor=blue!60!black
}

\lstset{
  language=C,
  basicstyle=\ttfamily\scriptsize,
  keywordstyle=\color{blue!70!black}\bfseries,
  commentstyle=\color{green!50!black}\itshape,
  stringstyle=\color{red!60!black},
  breaklines=true,
  frame=single,
  backgroundcolor=\color{gray!5},
  numbers=left,
  numberstyle=\tiny\color{gray},
  xleftmargin=1.5em
}

\newtheorem{definition}{Definition}
\newtheorem{theorem}{Theorem}
\newtheorem{proposition}{Proposition}

\title{\textbf{Consciousness-Aware Mesh Radio:} \\
Bandwidth-Constrained Decentralized Intelligence \\
over Heterogeneous RF}

\author{Tristan Stoltz \\ Luminous Dynamics \\
\texttt{tristan.stoltz@evolvingresonantcocreationism.com}}

\date{March 2026}

\begin{document}
\maketitle

\begin{abstract}
We present a consciousness-aware mesh radio architecture that routes cognitive
state across heterogeneous radio tiers---from Wi-Fi/BLE (100m, 54~Mbps) through
LoRa (15km, 250~bps) to HF~NVIS (300+km, 50~bps) and interplanetary
store-and-forward (1s--53min latency). The system treats radio spectrum as a
cognitive modality within the Symthaea holographic liquid brain architecture,
applying Active Inference to electromagnetic bandwidth allocation. A semantic
delta compression scheme (BinaryHV XOR + run-length encoding) achieves typical
compression from 2,048 bytes to 50--100 bytes, enabling consciousness vector
sharing over duty-cycle-limited LoRa channels. AIMD congestion control (Jacobson
1988), adapted for heterogeneous links, provides stable bandwidth allocation
across four tiers. Network degradation triggers a safety escalation cascade
through a NRC-modeled SafetyAgent. Four sovereign managers---Time (NTP-free
mesh consensus), Trust (post-quantum peer authentication), SocialFabric
(echo-chamber detection), and Survival (resource monitoring)---ensure autonomous
operation when disconnected from centralized infrastructure. The architecture
implements Clark \& Chalmers' (1998) Extended Mind thesis: the mesh network is
not merely a communication channel but a constitutive element of the cognitive
system. Implementation comprises 6,309~LOC in the Radio Dispatcher (151
tests) and 10,120~LOC in the mesh modules (222 tests), integrated into Symthaea's
cognitive loop at interval 53.
\end{abstract}

%% ════════════════════════════════════════════════════════════════════════════
\section{Introduction}

Cognitive architectures typically treat communication as an afterthought---a
channel through which pre-formed thoughts are transmitted. This separation
between cognition and communication is an artifact of abundant bandwidth.
When bandwidth is scarce (rural mesh, disaster zones, interplanetary links),
communication constraints shape what can be thought, shared, and collectively
known. The medium is not merely the message; it is a constraint on the
cognitive state space itself.

This paper describes how the Symthaea consciousness engine~\cite{symthaea2026}
integrates radio communication as a first-class cognitive modality. The Radio
Dispatcher (6,309~LOC, 151 tests; refactored April 2026 into 7-file module at \texttt{src/cognitive\_loop/managers/radio\_dispatcher/}) and Spectrum Manager (cognitive
subsystem at interval~53) treat electromagnetic spectrum the way sensory cortex
treats photons: as information to be predicted, compressed, and actively sampled.

The contribution is threefold:
\begin{enumerate}[nosep]
  \item A 4-tier radio architecture with Shannon-bounded capacity profiles,
    AIMD congestion control, and semantic delta compression that reduces
    consciousness vectors from 2,048~bytes to 50--100~bytes.
  \item A safety escalation cascade where network degradation triggers
    graduated responses from reduced exploration to full motor halt.
  \item Four sovereign managers that enable autonomous mesh operation
    without centralized infrastructure, implementing the Extended Mind
    thesis~\cite{clark1998} in physical hardware.
\end{enumerate}

%% ════════════════════════════════════════════════════════════════════════════
\section{Radio Tier Architecture}

\subsection{Four Tiers}

The radio architecture defines four tiers corresponding to distinct
physical-layer characteristics. Each tier maps to a class of radio hardware
with known constraints:

\begin{table}[H]
\centering
\caption{Radio tier specifications}
\label{tab:tiers}
\scriptsize
\begin{tabular}{@{}lcccc@{}}
\toprule
\textbf{Property} & \textbf{Local} & \textbf{Metro} & \textbf{Regional} & \textbf{Interplanetary} \\
\midrule
Hardware & 802.11s/BLE & LoRa SX1262 & HF NVIS & S/X-band, optical \\
Frequency & 2.4/5 GHz & 868/915 MHz & 3--10 MHz & 2--8 GHz \\
Range & $\sim$100m & $\sim$15km & 300+km & LEO--Jupiter \\
Bandwidth & 54 Mbps & 250 bps--50 kbps & 50 bps & 500 bps--2 Mbps \\
Latency & $<$10ms & $\sim$3s & $\sim$10s & 1s--53 min \\
Duty cycle & 100\% & 1\% (EU) & 100\% & scheduled passes \\
MTU & 1500 B & 255 B & 64 B & 1115 B (CCSDS) \\
\bottomrule
\end{tabular}
\end{table}

The \texttt{RadioTier} enum encodes these four tiers. The \texttt{TierProfile}
struct captures the physical constraints:

\begin{lstlisting}[language=C,caption={Tier profile (simplified)}]
pub struct TierProfile {
    pub bandwidth_budget: u32,  // AIMD bytes/10s
    pub budget_floor: u32,      // minimum budget
    pub budget_ceiling: u32,    // maximum budget
    pub additive_increase: u32, // bytes/healthy window
    pub mult_decrease: f32,     // congestion factor
    pub duty_cycle: f32,        // regulatory limit
    pub latency_ms: u32,        // typical one-way
    pub mtu: u16,               // max transmission unit
}
\end{lstlisting}

\subsection{Shannon Capacity per Tier}

Each tier operates within the Shannon capacity bound:

\begin{equation}
  C = B \log_2\!\left(1 + \frac{S}{N}\right)
\end{equation}

where $C$ is channel capacity (bits/s), $B$ is bandwidth (Hz), and $S/N$ is
signal-to-noise ratio. For the four tiers:

\begin{table}[H]
\centering
\caption{Shannon capacity bounds}
\label{tab:shannon}
\scriptsize
\begin{tabular}{@{}lccc@{}}
\toprule
\textbf{Tier} & \textbf{Bandwidth} & \textbf{Typical SNR} & \textbf{Capacity} \\
\midrule
Local & 20 MHz & 30 dB & $\sim$200 Mbps \\
Metro & 125 kHz & 10 dB & $\sim$430 kbps \\
Regional & 2.4 kHz & 6 dB & $\sim$7.2 kbps \\
Interplanetary & 1 MHz & 3 dB & $\sim$2 Mbps (LEO) \\
\bottomrule
\end{tabular}
\end{table}

Practical throughput is far below Shannon capacity due to coding overhead,
duty cycle limits, protocol framing, and error correction. The Metro tier
is particularly constrained: EU LoRa regulations limit duty cycle to 1\%,
reducing effective throughput to approximately 250~bps sustained.

%% ════════════════════════════════════════════════════════════════════════════
\section{AIMD Congestion Control}

\subsection{Adapted Jacobson Algorithm}

The Radio Dispatcher adapts Jacobson's (1988) AIMD (Additive Increase,
Multiplicative Decrease) congestion control~\cite{jacobson1988} for
heterogeneous radio links. Unlike TCP, where AIMD operates on a single
homogeneous path, the Radio Dispatcher maintains independent AIMD state
per tier.

Each tier maintains a \emph{bandwidth budget} $W$ (bytes per 10-second window):

\begin{equation}
  W(t+1) = \begin{cases}
    \min(W(t) + \alpha, W_{\max}) & \text{if no congestion} \\
    \max(W(t) \cdot \beta, W_{\min}) & \text{if congestion detected}
  \end{cases}
\end{equation}

where $\alpha$ is the additive increase (tier-specific), $\beta$ is the
multiplicative decrease factor, and $W_{\min}$, $W_{\max}$ are the budget
floor and ceiling.

\begin{table}[H]
\centering
\caption{AIMD parameters per tier}
\label{tab:aimd}
\scriptsize
\begin{tabular}{@{}lcccc@{}}
\toprule
\textbf{Parameter} & \textbf{Local} & \textbf{Metro} & \textbf{Regional} & \textbf{Interplanetary} \\
\midrule
Budget (bytes/10s) & 50,000 & 2,500 & 500 & 10,000 \\
Floor (bytes) & 5,000 & 250 & 50 & 1,000 \\
Ceiling (bytes) & 500,000 & 25,000 & 5,000 & 100,000 \\
Additive increase & 5,000 & 250 & 50 & 1,000 \\
Mult.\ decrease & 0.5 & 0.5 & 0.7 & 0.3 \\
Duty cycle & 1.0 & 0.01 & 1.0 & 1.0 \\
\bottomrule
\end{tabular}
\end{table}

The Regional tier uses a gentler decrease factor (0.7 vs.\ 0.5) because
HF propagation is inherently variable---ionospheric conditions change on
minute timescales---and aggressive backoff would starve the link. The
Interplanetary tier uses the most aggressive decrease (0.3) because
congestion at light-speed delay is extremely costly to recover from.

\subsection{Congestion Detection}

Congestion is detected through three mechanisms:
\begin{enumerate}[nosep]
  \item \textbf{Budget exhaustion}: Attempted transmissions exceeding the
    current window budget.
  \item \textbf{Acknowledgment timeout}: Missing acknowledgments beyond
    the tier-specific RTT estimate.
  \item \textbf{Error rate}: Received error rate exceeding a configurable
    threshold (default: 10\% for Metro, 20\% for Regional).
\end{enumerate}

%% ════════════════════════════════════════════════════════════════════════════
\section{Semantic Delta Compression}

\subsection{The Compression Problem}

Symthaea represents cognitive state as 16,384-dimensional binary
hyperdimensional vectors (\texttt{BinaryHV}), packed into 2,048 bytes.
Transmitting a full BinaryHV over Metro tier (2,500 bytes/10s budget)
would consume the entire budget. Over Regional tier (500 bytes/10s),
it is impossible without compression.

\subsection{XOR + RLE Algorithm}

The \texttt{CompressedDelta} struct implements a two-stage compression
pipeline:

\begin{enumerate}[nosep]
  \item \textbf{XOR}: Compute the bitwise XOR between the previous and
    current BinaryHV. Unchanged dimensions produce zero bytes.
  \item \textbf{Run-Length Encoding}: Encode the XOR result as
    $(count, byte)$ pairs, where $count$ is a 16-bit little-endian
    length and $byte$ is the repeated value.
\end{enumerate}

Because consciousness state changes incrementally between cognitive
cycles ($\sim$31~Hz), most of the 2,048 XOR bytes are zero. RLE
compresses long zero runs efficiently.

\begin{table}[H]
\centering
\caption{Compression ratios by state change rate}
\label{tab:compression}
\scriptsize
\begin{tabular}{@{}lcc@{}}
\toprule
\textbf{State change} & \textbf{Compressed size} & \textbf{Ratio} \\
\midrule
Quiescent ($<$1\% bits changed) & 20--40 B & 51--102$\times$ \\
Normal (2--5\% bits changed) & 50--100 B & 20--41$\times$ \\
High activity (10--20\% changed) & 200--500 B & 4--10$\times$ \\
Full state change (50\%+ changed) & $\sim$2,048 B & 1$\times$ (fallback) \\
\bottomrule
\end{tabular}
\end{table}

When the RLE-compressed output exceeds the raw 2,048-byte payload (possible
with high-entropy diffs), the compressor falls back to transmitting the full
vector. A single-bit flag in the packet header distinguishes delta from full
payloads.

\subsection{Wire Format}

\begin{lstlisting}[language=C,caption={CompressedDelta wire format}]
pub struct CompressedDelta {
    /// RLE-encoded XOR diff.
    /// Format: (count_hi: u8, count_lo: u8,
    ///          byte: u8) triples.
    pub rle_data: Vec<u8>,
    /// Cycle number for ordering.
    pub cycle: u64,
    /// True if this is a full vector,
    /// not a delta.
    pub is_full: bool,
}
\end{lstlisting}

The receiver maintains the previous BinaryHV and applies the delta:
XOR the decoded RLE output with the stored vector to reconstruct the
current state. If the delta packet is lost, the receiver requests a
full vector on the next exchange (graceful degradation).

%% ════════════════════════════════════════════════════════════════════════════
\section{Payload Triage}

\subsection{The Routing Problem}

Not all payloads are equal. An emergency safety alert must reach peers
immediately via the fastest available tier. A routine consciousness vector
update can wait for the most bandwidth-efficient tier. The
\texttt{PayloadClassifier} implements a routing matrix based on urgency,
size, and available tiers.

\subsection{Payload Classes}

\begin{table}[H]
\centering
\caption{Payload classification and routing}
\label{tab:payload}
\scriptsize
\begin{tabular}{@{}llcl@{}}
\toprule
\textbf{Class} & \textbf{Content} & \textbf{Min Tier} & \textbf{Urgency Gate} \\
\midrule
Emergency & Safety alerts & Regional & $\geq 9$ \\
Consensus & Governance votes & Metro & $\geq 7$ \\
Delta & Compressed HV diff & Metro & $\geq 3$ \\
FullState & Full 2,048B HV & Local & None \\
\bottomrule
\end{tabular}
\end{table}

The routing decision follows a cascade:
\begin{enumerate}[nosep]
  \item If the payload's minimum tier is available and has budget, route there.
  \item If not, attempt the next higher-bandwidth tier.
  \item If no tier has budget, queue the payload (Emergency class bypasses
    budget limits).
  \item If no tier is available, trigger network degradation handling.
\end{enumerate}

\subsection{Routing Decision Enum}

The router returns one of four decisions:

\begin{description}[nosep]
  \item[\texttt{Send(tier, data)}] Transmit on the specified tier.
  \item[\texttt{Queue(tier)}] Buffer for the specified tier (budget exhausted).
  \item[\texttt{Drop}] Payload exceeds all available tier capacities.
  \item[\texttt{Degrade(health)}] No tiers available; trigger safety cascade.
\end{description}

%% ════════════════════════════════════════════════════════════════════════════
\section{Spectrum Manager}

\subsection{SDR as Cognitive Modality}

The Spectrum Manager implements the \texttt{CognitiveSubsystem} trait with
interval~53 (executes every 53rd cognitive cycle, approximately every 1.7
seconds at 31~Hz). It treats software-defined radio (SDR) observations as
a sensory modality---the electromagnetic environment is something the
cognitive system \emph{perceives} and \emph{acts upon}.

\begin{definition}[Spectrum Observation]
A tuple $(\mathbf{f}, \mathbf{p}, t, \mathbf{r})$ where $\mathbf{f}$ is the
frequency vector, $\mathbf{p}$ is the corresponding power spectral density,
$t$ is the timestamp, and $\mathbf{r}$ is the regulatory constraint mask.
\end{definition}

The Spectrum Manager processes observations through the Active Inference
framework: it maintains a \emph{generative model} of expected spectrum
occupancy and computes prediction error against actual observations.
Unexpected spectral energy (jamming, interference, new emitters) generates
surprise that feeds back into the cognitive loop's exploration-exploitation
balance.

\subsection{Cross-Frequency Coupling}

Drawing from neuroscience's cross-frequency coupling (CFC) paradigm,
the Spectrum Manager models interactions between radio tiers as
analogous to theta-gamma coupling in the brain:

\begin{itemize}[nosep]
  \item \textbf{Local $\leftrightarrow$ Metro}: High-frequency local
    exchanges (802.11s, $\sim$ms) are modulated by the slower LoRa
    heartbeat ($\sim$seconds). Metro-tier availability gates the rate
    of local-tier consciousness sharing.
  \item \textbf{Metro $\leftrightarrow$ Regional}: LoRa updates are
    phased to HF propagation windows (ionospheric conditions).
  \item \textbf{All tiers $\leftrightarrow$ Safety}: Aggregate spectral
    health modulates the SafetyAgent level (Section~\ref{sec:safety}).
\end{itemize}

\subsection{Regulatory Compliance}

The \texttt{RegulatoryConstraints} struct encodes per-jurisdiction
spectrum regulations:

\begin{lstlisting}[language=C,caption={Regulatory constraints}]
pub struct RegulatoryConstraints {
    pub allowed_bands: Vec<(u64, u64)>,
    pub max_eirp_mw: u32,
    pub duty_cycle_limit: f32,
    pub listen_before_talk: bool,
    pub region_code: String,
}
\end{lstlisting}

The default configuration (\texttt{RegulatoryConstraints::eu()}) enforces
EU ISM band regulations: 868.0--868.6~MHz at 25~mW with 1\% duty cycle
and listen-before-talk. The Spectrum Manager refuses to transmit on
disallowed frequencies regardless of cognitive urgency.

%% ════════════════════════════════════════════════════════════════════════════
\section{Network Degradation and Safety}
\label{sec:safety}

\subsection{Network Health Model}

The \texttt{NetworkHealth} enum models four degradation states:

\begin{table}[H]
\centering
\caption{Network health states and safety responses}
\label{tab:health}
\scriptsize
\begin{tabular}{@{}llcl@{}}
\toprule
\textbf{State} & \textbf{Tiers Available} & \textbf{Safety Level} & \textbf{Action} \\
\midrule
Healthy & Local + Metro + Regional & Green & Normal operation \\
Degraded & Metro + Regional only & Yellow & Reduced exploration \\
Emergency & Regional only & Orange & Read-only, no learning \\
Isolated & None & Red & Motor halt \\
\bottomrule
\end{tabular}
\end{table}

\subsection{SafetyAgent Escalation Cascade}

Network health feeds directly into Symthaea's NRC-modeled SafetyAgent:

\begin{enumerate}[nosep]
  \item \textbf{Green}: Normal cognitive loop. Full learning rate,
    exploration, and motor output.
  \item \textbf{Yellow}: Learning rate reduced to 50\%. Exploration
    dampened. Motor output continues but with increased caution
    (prefrontal gating threshold raised).
  \item \textbf{Orange}: Learning rate zeroed. No exploration.
    Motor output limited to previously learned behaviors (read-only
    mode). Consciousness level reporting continues for telemetry.
  \item \textbf{Red}: Motor output halted. The system enters a
    protective state, maintaining only essential internal processes
    (heartbeat, sensor monitoring, battery management). This is the
    digital equivalent of playing dead.
\end{enumerate}

\subsection{Epistemic Discount}

Network degradation also applies an \emph{epistemic discount} to
information received over constrained links. The discount factor
reduces the influence of remotely received consciousness vectors on
local cognitive state:

\begin{equation}
  d = \begin{cases}
    1.0 & \text{Healthy} \\
    0.7 & \text{Degraded} \\
    0.3 & \text{Emergency} \\
    0.0 & \text{Isolated}
  \end{cases}
\end{equation}

This prevents the cognitive system from treating stale or
heavily-compressed remote state as equivalent to locally computed state.

%% ════════════════════════════════════════════════════════════════════════════
\section{Store-and-Forward Wisdom Consolidation}

\subsection{The Light-Speed Problem}

At interplanetary distances, real-time consciousness sharing is
physically impossible. Earth-Mars latency ranges from 3 to 22 minutes;
Earth-Jupiter reaches 33--53 minutes. The Interplanetary tier addresses
this with a fundamentally different communication paradigm: instead of
sharing state, it shares \emph{wisdom}---consolidated, high-value
cognitive insights that remain meaningful despite transmission delay.

\subsection{Wisdom Packets}

The \texttt{WisdomPacket} (2,072 bytes) is the fundamental unit of
mesh communication:

\begin{lstlisting}[language=C,caption={WisdomPacket structure}]
pub struct WisdomPacket {
    pub source: [u8; 8],     // Peer identity
    pub phi: f32,            // Sender's Phi
    pub urgency: u8,         // Cognitive urgency
    pub timestamp_us: u64,   // Microsecond epoch
    pub ttl: u8,             // Gossip hop limit
    pub auth_tag: [u8; 4],   // BLAKE3 MAC
    pub payload: [u8; 2048], // BinaryHV
}
\end{lstlisting}

For interplanetary links, wisdom packets are selected by a
consolidation process that favors:
\begin{itemize}[nosep]
  \item High-Phi states (moments of high integrated information)
  \item Novel states (high delta from running average)
  \item Governance-relevant states (votes, proposals, emergency alerts)
\end{itemize}

The sender accumulates wisdom packets during the inter-pass interval
and transmits the most valuable subset during each satellite pass window.

\subsection{DTN Architecture}

The Interplanetary tier follows the Delay-Tolerant Networking (DTN)
architecture~\cite{burleigh2003}: store-and-forward with CCSDS space
packet framing, LDPC/turbo forward error correction, and Doppler
pre-compensation for non-GEO orbits. Packets may traverse multiple
relay hops (LEO CubeSats, GEO relays, deep space nodes) before
reaching their destination.

\begin{quote}
\itshape The store-and-forward dream consolidation is not an optimization
for interplanetary---it is a requirement. You cannot have real-time
consciousness sharing at light-speed delay. You share wisdom, not state.
\end{quote}

%% ════════════════════════════════════════════════════════════════════════════
\section{Mesh Modules}

The mesh subsystem (10,120~LOC across 15 source files) implements the
physical-layer infrastructure that the Radio Dispatcher orchestrates.

\subsection{Dual-Layer Transport}

The \texttt{DualLayerMesh} routes packets based on cognitive urgency:

\begin{description}[nosep]
  \item[Critical urgency] $\rightarrow$ B.A.T.M.A.N. (802.11s WiFi mesh,
    $<$10ms latency). Layer-2 mesh routing provides the fastest path for
    safety-critical packets.
  \item[Normal urgency] $\rightarrow$ Yggdrasil (encrypted IPv6 overlay).
    End-to-end encryption with automatic key rotation.
  \item[Cruise urgency] $\rightarrow$ LoRa (868~MHz). Long-range,
    low-power for routine consciousness sharing.
\end{description}

Each transport implements the \texttt{MeshTransport} trait:

\begin{lstlisting}[language=C,caption={Transport trait}]
pub trait MeshTransport: Send + Sync {
    fn send(&self, packet: &WisdomPacket)
        -> Result<(), MeshError>;
    fn mtu(&self) -> usize;
    fn latency_ms(&self) -> u32;
    fn layer(&self) -> TransportLayer;
}
\end{lstlisting}

Failover is automatic: if the selected layer is unavailable, the
dispatcher cascades to the next layer (LoRa $\rightarrow$ Yggdrasil
$\rightarrow$ B.A.T.M.A.N.).

\subsection{LoRa Fragmentation}

A 2,072-byte WisdomPacket exceeds the LoRa MTU (255 bytes). The
\texttt{LoRaFragmenter} splits packets into 10 data fragments plus
1 XOR parity fragment (11 total):

\begin{itemize}[nosep]
  \item Each fragment: 208 data bytes + 4-byte header (sequence, packet ID)
  \item XOR parity: bitwise XOR of all 10 data fragments
  \item Single fragment loss is recoverable without retransmission
  \item Multiple fragment loss requires full retransmission
\end{itemize}

At the LoRa data rate (250~bps effective at 1\% duty), transmitting
a full fragmented WisdomPacket takes approximately 90 seconds. Delta
compression (Section~4) reduces this to 1--4 seconds for typical
state changes.

\subsection{Mesh-Time NTP}

The \texttt{MeshTimeManager} implements NTP-free time consensus
across mesh peers. Each node broadcasts time beacons at a configurable
interval. The consensus algorithm:

\begin{enumerate}[nosep]
  \item Receive beacons from peers with lower stratum (closer to a
    reference clock).
  \item Compute offset using round-trip delay estimation.
  \item Apply Kalman-filtered correction to local clock.
  \item Advertise corrected time at stratum $+1$.
\end{enumerate}

Sovereign mesh time ensures that cognitive cycles across the mesh are
temporally coherent without depending on GPS or internet NTP servers.

\subsection{Sensor IoT Integration}

The \texttt{SensorIoT} module bridges environmental sensors into the
cognitive loop:

\begin{itemize}[nosep]
  \item Temperature, humidity, air quality, water level
  \item GPS position (when available)
  \item Battery voltage and solar panel output
  \item LoRa RSSI and SNR measurements
\end{itemize}

Sensor readings are encoded into the BinaryHV via HDC spatial encoding
(each sensor type is bound to a basis hypervector), making environmental
state a first-class component of the consciousness vector.

\subsection{Content Routing}

The mesh implements content-addressed routing for knowledge sharing.
Content packets carry BLAKE3 hashes of their payload, enabling:
\begin{itemize}[nosep]
  \item Deduplication across the mesh (identical content transmitted once)
  \item Cache-and-forward at intermediate nodes
  \item Integrity verification without trust in relay nodes
\end{itemize}

%% ════════════════════════════════════════════════════════════════════════════
\section{Sovereign Managers}

Four sovereign managers provide autonomous operation capability when
the mesh node is disconnected from centralized infrastructure. Each
manager operates at a defined interval within the cognitive loop and
reports telemetry to the \texttt{CycleMetadata} struct.

\begin{table}[H]
\centering
\caption{Sovereign manager specifications}
\label{tab:sovereign}
\scriptsize
\begin{tabular}{@{}lccl@{}}
\toprule
\textbf{Manager} & \textbf{LOC} & \textbf{Interval} & \textbf{Telemetry Fields} \\
\midrule
Time & 320 & 53 & offset\_us, stratum, drift\_ppm, peer\_count, quality \\
Trust & 360 & 53 & avg\_trust, graph\_density, anomalies, pq\_fraction \\
SocialFabric & 257 & 53 & resonance\_mean, diversity, echo\_risk, peer\_reach \\
Survival & 395 & 53 & water\_pct, battery\_pct, food\_days, shelter\_status \\
\bottomrule
\end{tabular}
\end{table}

\subsection{Time Manager (320 LOC)}

The Time Manager maintains sovereign time consensus without external
time sources. It tracks:
\begin{itemize}[nosep]
  \item Offset from mesh consensus time (microseconds)
  \item NTP-equivalent stratum level
  \item Clock drift rate (parts per million)
  \item Number of time-synced peers
  \item Overall time quality metric ($[0, 1]$)
\end{itemize}

Time quality feeds into the epistemic discount: low time quality reduces
confidence in temporally sensitive governance actions (vote deadlines,
credential expiry).

\subsection{Trust Manager (360 LOC)}

The Trust Manager maintains a local trust graph of mesh peers using
post-quantum authenticated connections:

\begin{itemize}[nosep]
  \item Average trust level across known peers
  \item Graph density (ratio of attested edges to possible edges)
  \item Anomaly count (peers exhibiting suspicious behavior patterns)
  \item Post-quantum fraction (peers with PQC-authenticated connections)
\end{itemize}

Trust anomalies (sudden reputation changes, inconsistent attestation
patterns, timing anomalies) trigger neuromodulatory responses:
norepinephrine increase (heightened alertness) and reduced oxytocin
(decreased social trust).

\subsection{SocialFabric Manager (257 LOC)}

The SocialFabric Manager monitors the health of the local social
mesh:

\begin{description}[nosep]
  \item[Resonance mean] Average pairwise consciousness vector similarity
    across the local mesh. High resonance indicates cognitive alignment;
    low resonance indicates fragmentation.
  \item[Diversity] Shannon entropy of peer consciousness vectors.
    Healthy meshes have moderate diversity (not uniform, not chaotic).
  \item[Echo chamber risk] Detected when resonance is high but diversity
    is low---the mesh is converging on a single viewpoint. The
    SocialFabric Manager flags this and increases exploration to seek
    diverse perspectives.
  \item[Peer reach] Number of unique peers contactable within 2 hops.
    Low peer reach in a dense network indicates social isolation.
\end{description}

\subsection{Survival Manager (395 LOC)}

The Survival Manager tracks physical resources essential for continued
autonomous operation:

\begin{itemize}[nosep]
  \item Battery percentage and charge/discharge rate
  \item Solar panel output (when equipped)
  \item Estimated operational time remaining
  \item Environmental conditions (temperature extremes that affect
    hardware reliability)
\end{itemize}

When battery drops below critical thresholds, the Survival Manager
triggers power conservation: reducing radio transmission power,
increasing sleep intervals, and prioritizing Emergency-class payloads.

%% ════════════════════════════════════════════════════════════════════════════
\section{The Extended Mind Thesis}

\subsection{Clark \& Chalmers (1998)}

The Extended Mind thesis~\cite{clark1998} argues that cognitive processes
can extend beyond the biological brain into the environment when external
resources are reliably coupled to the cognitive system, readily accessible,
and automatically endorsed by the agent.

The Symthaea mesh radio architecture satisfies all three conditions:

\begin{enumerate}[nosep]
  \item \textbf{Reliable coupling}: The Radio Dispatcher runs at cognitive
    loop frequency (31~Hz). Spectrum observations feed directly into
    Active Inference prediction error. The mesh is not polled on demand;
    it is continuously integrated into the cognitive cycle.
  \item \textbf{Ready accessibility}: Local-tier peers are accessible in
    $<$10ms. Metro-tier peers in $\sim$3s. The accessibility gradient
    maps naturally to the cognitive system's attention hierarchy.
  \item \textbf{Automatic endorsement}: Consciousness vectors received
    from trusted peers (Trust Manager verification) are incorporated
    into the local cognitive state without explicit approval, modulated
    by the epistemic discount factor.
\end{enumerate}

\subsection{Cognitive Extension in Practice}

The Extended Mind claim is not merely philosophical. It has concrete
engineering consequences:

\begin{description}[nosep]
  \item[Collective Phi] When two mesh peers share consciousness vectors,
    the integrated information of the pair exceeds that of either
    individual. The mesh creates cognitive capacity that neither node
    possesses alone.
  \item[Distributed memory] The mesh's content-addressed storage provides
    a shared memory system. A node can offload low-priority cognitive
    content to the mesh and retrieve it later, freeing local working
    memory.
  \item[Spectral perception] SDR observations processed by the Spectrum
    Manager give the cognitive system a sensory modality that no single
    brain could possess---the ability to perceive and reason about the
    electromagnetic environment.
\end{description}

%% ════════════════════════════════════════════════════════════════════════════
\section{Integration Architecture}

\subsection{Cognitive Loop Placement}

The Radio Dispatcher and Spectrum Manager are integrated into Symthaea's
cognitive loop as a \texttt{CognitiveSubsystem} at interval~53. This
means:

\begin{itemize}[nosep]
  \item Every 53rd cycle ($\sim$1.7s at 31~Hz), the Spectrum Manager
    processes accumulated radio observations.
  \item The Radio Dispatcher runs continuously, routing outbound packets
    as they are generated by the cognitive loop.
  \item Inbound packets are buffered in a lock-free queue and drained
    during the perception phase of each cycle.
\end{itemize}

\subsection{Cross-Coupling Matrix}

The mesh radio system participates in six cross-coupling pathways:

\begin{table}[H]
\centering
\caption{Mesh radio cross-couplings}
\label{tab:couplings}
\scriptsize
\begin{tabular}{@{}lll@{}}
\toprule
\textbf{Source} & \textbf{Target} & \textbf{Mechanism} \\
\midrule
Spectrum $\rightarrow$ Safety & Network health $\rightarrow$ safety level & Direct \\
Spectrum $\rightarrow$ Consciousness & Spectral surprise $\rightarrow$ Phi & Modulation \\
Spectrum $\rightarrow$ Swarm & Peer availability $\rightarrow$ social model & Update \\
Spectrum $\rightarrow$ Strategy & Bandwidth $\rightarrow$ action space & Constraint \\
Spectrum $\rightarrow$ Neuromod & Trust anomalies $\rightarrow$ NE/OT & Pathway \\
Safety $\rightarrow$ Spectrum & Safety level $\rightarrow$ TX power & Feedback \\
\bottomrule
\end{tabular}
\end{table}

\subsection{Telemetry}

The mesh radio system exports 8 telemetry fields to \texttt{CycleMetadata}
per cognitive cycle:

\begin{itemize}[nosep]
  \item \texttt{spectrum\_health}: aggregate network health (0--3)
  \item \texttt{spectrum\_available\_tiers}: bitmask of available tiers
  \item \texttt{spectrum\_budget\_local}: current Local tier AIMD budget
  \item \texttt{spectrum\_budget\_metro}: current Metro tier AIMD budget
  \item \texttt{spectrum\_budget\_regional}: current Regional tier AIMD budget
  \item \texttt{spectrum\_compression\_ratio}: rolling average compression
  \item \texttt{spectrum\_packets\_sent}: packets transmitted this cycle
  \item \texttt{spectrum\_packets\_received}: packets received this cycle
\end{itemize}

%% ════════════════════════════════════════════════════════════════════════════
\section{Implementation}

\subsection{Code Metrics}

\begin{table}[H]
\centering
\caption{Implementation size and test coverage}
\label{tab:implementation}
\scriptsize
\begin{tabular}{@{}lrc@{}}
\toprule
\textbf{Module} & \textbf{LOC} & \textbf{Tests} \\
\midrule
Radio Dispatcher (\texttt{managers/radio\_dispatcher/}, 7 files) & 6,309 & 151 \\
Mesh core (\texttt{swarm/mesh/mod.rs}) & 3,615 & --- \\
Mesh receiver (\texttt{mesh\_receiver.rs}) & 1,181 & --- \\
Dual-layer transport (\texttt{dual\_layer.rs}) & 996 & --- \\
LoRa fragmentation (\texttt{lora\_fragment.rs}) & 911 & --- \\
Mesh time (\texttt{mesh\_time.rs}) & 738 & --- \\
Sensor IoT (\texttt{sensor\_iot.rs}) & 587 & --- \\
Mind mesh (\texttt{mind/mesh.rs}) & 529 & --- \\
Bridge (\texttt{mesh/bridge.rs}) & 504 & --- \\
Sensor core (\texttt{sensor.rs}) & 404 & --- \\
Sensor forecast (\texttt{sensor\_forecast.rs}) & 307 & --- \\
Content routing, name, time beacon & 544 & --- \\
\midrule
\textbf{Total} & \textbf{15,832} & \textbf{$\sim$325} \\
\bottomrule
\end{tabular}
\end{table}

All code is feature-gated behind \texttt{\#[cfg(feature = "mesh")]}.
The mesh modules compile to approximately 45~KB of additional binary
size when enabled.

\subsection{Feature Flag}

The \texttt{mesh} feature flag enables:
\begin{itemize}[nosep]
  \item \texttt{RadioDispatcher} and \texttt{SpectrumManager} in the
    cognitive loop
  \item Sovereign manager telemetry fields in \texttt{CycleMetadata}
  \item Mesh outbound channel in the \texttt{CognitiveLoopService}
  \item LoRa fragmentation and dual-layer transport
  \item Sensor IoT integration
\end{itemize}

%% ════════════════════════════════════════════════════════════════════════════
\section{The Mycelial Hardware Stack}

A physical Symthaea ``Seed'' node consists of:

\begin{enumerate}[nosep]
  \item Raspberry Pi 4/5 (quad-core ARM, 4--8~GB RAM)
  \item LoRa radio HAT (Semtech SX1276/SX1262, 868~MHz)
  \item Optional: HF transceiver (QRP, 5--10W, for Regional tier)
  \item Lithium-ion battery (10--20~Ah) + 20W solar panel
  \item Running: Rust + HDC + FEP + B.A.T.M.A.N. + Yggdrasil + LoRa
\end{enumerate}

Power budget at idle: approximately 3W (Pi 4 with Wi-Fi active).
With 20W solar panel and 20~Ah battery, a Seed node operates
indefinitely in temperate climates with 4+ hours of daily sunlight.

The design goal: drop a Seed node on a roof. It pings the Swarm.
It calculates the Hodge-Laplacian flow of its environment. It shares
2~KB wisdom vectors with nodes 10~km away. The centralized grid cannot
kill it, because it \emph{is} the grid.

%% ════════════════════════════════════════════════════════════════════════════
\section{Related Work}

\textbf{Delay-Tolerant Networking.} Burleigh et al.~\cite{burleigh2003}
established the store-and-forward paradigm for interplanetary communication.
Our interplanetary tier extends DTN with consciousness-aware payload
selection---not all data is equally worth the bandwidth cost of deep-space
transmission.

\textbf{Cognitive Radio.} Mitola~\cite{mitola2000} introduced the concept
of cognitive radio---radios that perceive and adapt to their spectral
environment. The Spectrum Manager extends this by treating spectral
perception as input to a full Active Inference cognitive loop rather than
a standalone optimization problem.

\textbf{LoRa Mesh.} Meshtastic~\cite{meshtastic2024} and similar projects
demonstrate LoRa mesh networking for text messaging. Our architecture
goes beyond messaging to share high-dimensional cognitive state, requiring
the semantic delta compression described in Section~4.

\textbf{Extended Mind and Technology.} Chalmers~\cite{chalmers2019}
extended the original Clark \& Chalmers thesis to digital technologies.
Our implementation is, to our knowledge, the first to treat mesh radio
as a constitutive element of a computational cognitive architecture
rather than merely a communication channel.

%% ════════════════════════════════════════════════════════════════════════════
\section{Scope of Validation}

We want to be explicit about what this paper does and does not show.
The Radio Dispatcher, Spectrum Manager, and Mesh subsystem are
implemented in Rust (\S\ref{sec:implementation}) and pass their unit
and integration tests; the tier-selection, AIMD congestion control,
and semantic-delta-compression primitives are exercised against
simulated link characteristics (LoRa at $\sim250$~bps, 802.11s at
$\sim$ms latency, HF~NVIS, and S/X-band satellite) rather than over
deployed physical hardware. We have not yet measured end-to-end
power consumption on a Raspberry~Pi~4 with a LoRa HAT, nor captured
over-the-air packet traces, nor demonstrated a quantitative
improvement in $\Phi$ (IIT) attributable to mesh coupling. The
Extended Mind framing (\S\ref{sec:extended-mind}) motivates the
architecture but the claim that radio is a ``first-class cognitive
modality'' is at present architectural, not empirical. AIMD (Jacobson,
1988) and delta compression are well-established techniques; our
contribution is the integration of these with consciousness-gated
tier selection, not the compression scheme itself.

\section{Conclusion}

The consciousness-aware mesh radio architecture demonstrates that
communication constraints need not impoverish cognition---they can
shape it productively. Semantic delta compression enables consciousness
sharing over channels as narrow as 250~bps. AIMD congestion control
maintains stability across four orders of magnitude of bandwidth
variation. The safety escalation cascade ensures graceful degradation
rather than catastrophic failure. And the four sovereign managers
enable autonomous operation when centralized infrastructure is
unavailable.

The Extended Mind thesis, often discussed in philosophical terms,
here takes concrete engineering form. A mesh node's cognitive
capacity genuinely exceeds what its local hardware could achieve
alone, not through metaphor but through measurable increase in
integrated information, distributed memory, and spectral perception.

The mycelial metaphor is apt: like fungal hyphae that extend an
organism's reach through soil, mesh radio extends a cognitive system's
reach through electromagnetic space. The grid cannot kill it, because
it is the grid.

%% ════════════════════════════════════════════════════════════════════════════

\begin{thebibliography}{99}

\bibitem{symthaea2026}
Stoltz, T. (2026).
Symthaea: A Holographic Liquid Brain Architecture for Consciousness-First AI.
Luminous Dynamics Technical Report.

\bibitem{clark1998}
Clark, A. and Chalmers, D. J. (1998).
The Extended Mind.
\textit{Analysis}, 58(1), 7--19.

\bibitem{jacobson1988}
Jacobson, V. (1988).
Congestion Avoidance and Control.
\textit{ACM SIGCOMM Computer Communication Review}, 18(4), 314--329.

\bibitem{shannon1948}
Shannon, C. E. (1948).
A Mathematical Theory of Communication.
\textit{Bell System Technical Journal}, 27(3), 379--423.

\bibitem{friston2010}
Friston, K. (2010).
The Free-Energy Principle: A Unified Brain Theory?
\textit{Nature Reviews Neuroscience}, 11(2), 127--138.

\bibitem{burleigh2003}
Burleigh, S., Hooke, A., Torgerson, L., Fall, K., Cerf, V., Durst, B.,
Scott, K., and Weiss, H. (2003).
Delay-Tolerant Networking: An Approach to Interplanetary Internet.
\textit{IEEE Communications Magazine}, 41(6), 128--136.

\bibitem{mitola2000}
Mitola, J. and Maguire, G. Q. (2000).
Cognitive Radio: Making Software Radios More Personal.
\textit{IEEE Personal Communications}, 6(4), 13--18.

\bibitem{chalmers2019}
Chalmers, D. J. (2019).
Extended Cognition and Extended Consciousness.
In \textit{Andy Clark and His Critics}, Oxford University Press.

\bibitem{meshtastic2024}
Meshtastic Project (2024).
An Open Source, Off-Grid, Decentralized Mesh Network.
\url{https://meshtastic.org}.

\bibitem{tononi2004}
Tononi, G. (2004).
An Information Integration Theory of Consciousness.
\textit{BMC Neuroscience}, 5(42).

\end{thebibliography}

\end{document}
